\section{Một câu hỏi không cần thí nghiệm}
\label{sec:ch13-cau-hoi-khong-can-thi-nghiem}

\index{giới hạn lý thuyết!loại bằng chứng|(}%
Bài báo mà chương này đọc là một bài lý thuyết thuần túy của Shrisha Rao, dài 65
trang, một tác giả duy nhất, và tính tới lúc viết chương này thì vẫn là bản tiền
ấn phẩm trên arXiv, không có nơi công bố nào ghi
kèm~\autocite{p26limits}. Nó chứng minh định lý và
không chạy gì cả. Ở phần kết, chính tác giả xếp việc kiểm chứng bằng thực
nghiệm vào phần công việc tương lai: kiểm chứng bằng thực nghiệm những dự đoán
lý thuyết của bài là công việc tương lai quan trọng, và bước đầu có thể là khớp
vài hàm đa thức đơn giản bằng cây quyết định với độ phức tạp thay đổi dần, rồi
đối chiếu với các chặn mà bài đưa ra~\autocite{p26limits}. Trong cả bài không có
một bảng kết quả nào, không một hình vẽ nào, và không một bộ dữ liệu nào.

Vì thế loại bằng chứng ở đây khác hẳn loại bằng chứng mà bốn chương vừa qua thu
thập, và chỗ khác nhau đáng nói ra trước khi vào nội dung. Một số đo nói về một
lần chạy: bài báo neo đo được AUROC quanh mức ngẫu nhiên trên tập BonaFide của
nó, và câu ấy đúng cho tám chỉ số ấy, trên tập ấy, với mười mô hình ấy. Muốn nó
đúng rộng hơn thì phải chạy thêm. Một chặn thì nói theo chiều ngược lại: nó phát
biểu rằng trong mọi trường hợp thỏa các giả thiết đã nêu, một sự việc không xảy
ra được. Muốn nó đúng rộng hơn thì không phải chạy gì thêm, mà phải nới các giả
thiết.

Cái giá phải trả cho sức mạnh ấy vì thế nằm trọn ở chỗ giả thiết. Một số đo có thể sai
vì đo hỏng, và người đọc phát hiện được bằng cách đo lại. Một chặn thì không sai
bao giờ nếu chứng minh đúng; điều nó có thể làm là ràng buộc một đối tượng không
phải đối tượng người đọc quan tâm. Cho nên với một kết quả loại này, câu hỏi
kiểm chứng không phải \enquote{có đúng không} mà là \enquote{đúng về cái gì}, và
câu trả lời nằm ở các định nghĩa chứ không ở phần chứng minh.

\index{giới hạn lý thuyết!đối tượng của bài báo}%
Với bài báo này, câu trả lời đọc được từ chỗ bài không nói. Rà toàn văn cả 65
trang, chữ \enquote{faithfulness} và chữ \enquote{faithful} không xuất hiện lần
nào; \enquote{ground truth} cũng vậy; \enquote{attribution} xuất hiện đúng một
lần, trong một câu liệt kê ở phần mở đầu. Không một họ chỉ số nào của
chương~\ref{ch:do-do-trung-thuc} có mặt, dưới bất kỳ tên nào.

Vậy đối tượng của bài là sai số giữa đầu ra của lời giải thích và đầu ra của mô
hình, dưới một hạn mức đặt lên độ dài mô tả của lời giải thích. Đó là độ khớp,
dòng đầu của bảng~\ref{tab:ch01-bon-quan-he}, và
mục~\ref{sec:ch10-nam-nhom-van-de} đã tách nó khỏi độ trung thực. Một cái trần
đặt lên đại lượng ấy không tự động là một cái trần đặt lên độ trung thực, và
mục~\ref{sec:ch13-tran-noi-gi} là chỗ chương này nói nó ràng buộc được gì và
không ràng buộc được gì. Đây là lần thứ năm Phần~IV gặp một bài báo có đối tượng
lệch khỏi đối tượng của phần, và là lần đầu tiên chỗ lệch nằm ở loại phát biểu
chứ không ở họ lời giải thích.

\index{lý thuyết thông tin thuật toán}%
Công cụ mà bài dùng là \tn{lý thuyết thông tin thuật toán}{algorithmic
information theory}, viết tắt AIT, một nhánh lý thuyết đo lượng thông tin trong
một đối tượng riêng lẻ bằng độ dài chương trình ngắn nhất sinh ra đối tượng ấy.
Chỗ này khác lý thuyết thông tin của Shannon, vốn đo lượng thông tin trung bình
của một nguồn ngẫu nhiên: AIT hỏi một chuỗi bit cụ thể nén được xuống bao nhiêu,
chứ không hỏi một phân phối có bao nhiêu entropy. Sự khác biệt ấy là lý do chọn
nó ở đây, vì một mô hình đã huấn luyện xong là một đối tượng cụ thể chứ không
phải một phân phối.

Còn một điều đáng ghi trước khi vào mô hình toán, và nó là quan sát của sách chứ
không phải phát biểu của bài. Bốn chương vừa qua đều mắc chung một khó khăn: mọi
kết luận của chúng đều phải đi qua một dụng cụ đo mà chính lĩnh vực chưa kiểm
định. Chương~\ref{ch:chi-so-khong-do-do-trung-thuc} phải dựng ground truth mới
nói được điều gì, chương~\ref{ch:con-nguoi-vang-mat} phải hỏi 167 người. Kết quả
của chương này không đi qua dụng cụ nào. Nó là kết quả đầu tiên trong Phần~IV
không cần một phép đo nào để phát biểu, và cũng chính vì thế mà nó không trả lời
được câu hỏi về dụng cụ đo.

\index{giới hạn lý thuyết!loại bằng chứng|)}%
